\input{preamble}

\begin{document}

\title{The Price of a Frozen Surrogate in Support-Aware Compression:\\
Separation, Starvation and the Cost of Repair}

\author{\IEEEauthorblockN{Vu Phan An Nguyen}
\IEEEauthorblockA{\textit{Rider University} \\
Lawrenceville, NJ, USA}}

\maketitle

\begin{abstract}
A compression is chosen by comparing losses. A learner that chooses one under noise faces an
information game whose alternatives are the worlds in which its choice would be wrong. That
alternative set is routinely replaced by a surrogate: the extrema realising each loss are
frozen at their current cells and the ordering constraints are dropped. We ask what this
costs, and what it costs to undo. The model is correct, the stopping test is valid and the
decision is unchanged; only the geometry used to score designs is frozen.

\emph{Separation after optimisation.} Take a construction with four extremum cells and $r$
pairs of middle cells at variance $M$, and let both games range over the same centred class.
The frozen game has value exactly $\frac{9}{512}$, while the true game lies between
$\frac{9}{128(4+2rM)}$ and $\frac{1}{8rM}$. The ratio is therefore $\Theta(\frac{1}{rM})$, and
the gap is present at $M=1$.

\emph{Starvation.} Every design that attains the surrogate's optimum puts zero weight on the
middle cells, and so collects no information against the alternatives the decision turns on.
A design within $\varepsilon$ of that optimum collects at most $\frac{64\varepsilon}{9rM}$.

\emph{The price of repair.} Giving up a fraction $\beta$ of the surrogate's value buys back
$\Theta(\frac{\beta}{rM})$ of true information, on the game the decision defines. Operationally,
every $\delta$-correct policy must give up a \emph{budget} of surrogate value. Let $N$ be the
expected number of queries and $\beta$ the relative sacrifice of the average design. Then
$N\beta\ge8rM\log(1/\delta)(1-o(1))$. The rate-optimal learner pays at most
$\frac{128(4+2rM)}{9}\log(1/\delta)(1+o(1))$, so the budget is $\Theta(rM\log(1/\delta))$. The
same bound gives a forcing law: a learner whose informative queries grow like $cT^{a}$ needs
$\mathbb{E}T\gtrsim\log^{1/a}(1/\delta)$.

\emph{Learners.} An ideal learner, defined and proved here, is $\delta$-correct and attains
$\mathbb{E}\tau_{\delta}/\log(1/\delta)\to1/\Gamma_{\rm true}$. The learner that tracks the
surrogate needs $\Omega(\log^{2}(1/\delta))$ samples, both when it tracks the surrogate's optimum
at the true world and when it re-solves the surrogate at its own estimate.

All theorems are machine-checked in Lean~4. Measurements on finite panels are reported as such.
\end{abstract}

\begin{IEEEkeywords}
compression, alternative sets, experimental design, surrogate objectives, sequential testing,
formal verification
\end{IEEEkeywords}

\section{Introduction}
\label{sec:intro}

\subsection{The question}

Let a world be $\theta$, let $T$ be a retained coordinate set, and let $L_{\theta,K}(T)$ be the
compression loss of $T$ on the deployment support $K$. A learner must return a $T$ whose
regret is at most $\varepsilon$, and it pays for observations. Its cost is governed by an
information game. The inner infimum of that game ranges over the \emph{alternatives}: the
worlds in which the returned $T$ would be wrong.

Computing that infimum requires the alternative set. The natural shortcut freezes the cells
realising each maximum and minimum, and asks how far the observations must move to reverse
one linear contrast. We call this the \emph{frozen surrogate}. The observation model stays
correct, the stopping test stays valid and the decision is unchanged. Only the geometry used
to score designs is altered. This paper measures what that alteration costs, and what it
costs to undo.

\subsection{Results}

Four results carry the paper.
\begin{enumerate}
\item[1)] \textbf{Separation after optimisation} (Theorem~\ref{thm:sep}). Both games optimise
the allocation and range over the same centred class. Even so, the true game is a factor
$\Theta(rM)$ below the frozen one, and the factor is already present at unit variance.
\item[2)] \textbf{Starvation} (Theorem~\ref{thm:starve}). Every optimum of the surrogate puts
zero weight on the cells that carry the decision. An approximate optimum leaks in proportion
to its suboptimality.
\item[3)] \textbf{The price of repair} (Theorems~\ref{thm:frontier}, \ref{thm:sacrifice},
\ref{thm:dilute}, \ref{thm:forcing}). A design that sacrifices a fraction $\beta$ of surrogate value recovers
$\Theta(\beta/rM)$ of true information. Operationally, a $\delta$-correct policy must sacrifice
a budget of $8rM\log(1/\delta)$ query-units of surrogate value to first order, and the
rate-optimal learner pays that budget up to a factor of at most $\frac{32}{9}+\frac{64}{9rM}$.
The budget can be spread over any run length: the relative sacrifice falls as $N_{0}/N$ while
the policy stays $\delta$-correct (Theorem~\ref{thm:dilute}).
A learner that meets its demand for informative queries only through forcing at rate
$t^{a-1}$ pays $\log^{1/a}(1/\delta)$.
\item[4)] \textbf{Learners} (Theorems~\ref{thm:ideal}, \ref{thm:surlearner}). A completely
defined ideal learner attains the true rate. The learner that tracks the surrogate, at the
truth or re-solved at its own estimate, pays a whole extra power of $\log(1/\delta)$.
\end{enumerate}
Three further facts are what make these statements about compression rather than about a
parameter class: the construction is an observation model, the alternative set is exactly a
union of polyhedra, and a margin certificate is not a decision (Section~\ref{sec:model}).

\subsection{Prior work, and what does not follow from it}

Information games with several correct answers, and alternative sets that decompose into
finitely many convex regions, are developed in the literature. Degenne and
Koolen~\cite{degenne2019} treat multiple correct answers, with a lower bound and an algorithm
that attains it; we inherit that machinery and do not claim it. Wang and You~\cite{wangyou2026}
study ranking-and-selection with multiple correct answers and non-answerable estimates. That
an information proxy can select a design lacking identifiability is
known~\cite{overstall2022}. So is the fact that acquisition should target the decision rather
than the predictor~\cite{filstroff2021}, error amplification under
misspecification~\cite{tang2026}, and decision-focused sequential design with a failing
prediction-focused instance~\cite{wan2026}. When a proxy suffices for a decision is studied
directly~\cite{schutte2026}. Bennouna, Bennouna, Amin and Ozdaglar~\cite{bennouna2026} study
data informativeness in linear optimisation under uncertainty. We do not extend their
statements to noisy compression, which they do not make.

The generic statements are therefore taken. What remains is narrower. The model is correct,
the test is valid and the decision is unchanged, and still the frozen geometry changes what a
design can learn and what a learner pays. Table~\ref{tab:prior} states the difference at the
level of theorems.

\begin{table}[h]
\caption{Nearest sources, and a theorem of this paper that does not follow from each.}
\label{tab:prior}
\centering
\footnotesize
\begin{tabular}{@{}p{0.30\columnwidth}p{0.62\columnwidth}@{}}
\toprule
Nearest source & Theorem here that does not follow \\
\midrule
Proxy insufficiency~\cite{schutte2026,overstall2022} & The bracket
$\frac{4}{4+2rM}\le\Gamma_{\rm true}/\Gamma_{\rm fr}\le\frac{64}{9rM}$ after optimisation, and
zero information at \emph{every} surrogate optimum (Thms.~\ref{thm:sep},
\ref{thm:starve}) \\
Multiple-answer track-and-stop~\cite{degenne2019} & The repair frontier
$\Theta(\beta/rM)$ and the sacrifice budget $N\beta\ge8rM\log(1/\delta)(1-o(1))$
(Thms.~\ref{thm:frontier}, \ref{thm:sacrifice}) \\
Decision-focused design~\cite{wan2026} & The $\Omega(\log^{2}(1/\delta))$ cost of the
learner that re-solves the surrogate at its estimate, against the ideal learner's
$\Theta(\log(1/\delta))$ (Thms.~\ref{thm:ideal}, \ref{thm:surlearner}) \\
\bottomrule
\end{tabular}
\end{table}

\subsection{What is proved, and how}

Three evidence classes appear, and they are never interchanged.
\begin{itemize}
\item \PROVED: a theorem stated at the paper's objects and checked in
Lean~4~\cite{demoura2021} against Mathlib~\cite{mathlib2020}.
\item \CERT: an exact rational certificate.
\item \NUM: a floating-point or Monte Carlo result.
\end{itemize}
Each theorem carries one line naming its Lean declaration. The supplement gives the proof
architecture, every supporting declaration, and the status of results superseded during
preparation. All modules compile with no \texttt{sorry} and no \texttt{native\_decide}, and
their axioms lie within \{\texttt{propext}, \texttt{Classical.choice}, \texttt{Quot.sound}\}.
Every hypothesis of a headline is either a model assumption displayed in it or discharged.

\section{The model, the decision and its alternatives}
\label{sec:model}

\paragraph{The loss} Parametrise the world linearly, $F_{\theta}(x)=\phi(x)^{\top}\theta$ on a
finite domain with a fixed reference and a deployment support $K$. For a retained set $T$ the
residual at a support point is a linear form $a_{T}(x)^{\top}\theta$, and
\begin{equation}
L_{\theta,K}(T)=\tfrac12\Bigl[\max_{x\in K}a_{T}(x)^{\top}\theta
-\min_{x\in K}a_{T}(x)^{\top}\theta\Bigr].
\label{eq:loss}
\end{equation}
Query $e$ returns $Y_{e}\sim\mathcal N(\theta_{e},\nu_{e})$, and a design is an allocation
$w\in\Delta$. For a set of alternatives $\mathcal A$, the value of a design and the game are
\begin{equation}
J(w)=\inf_{\lambda\in\mathcal A}\ \tfrac12\sum_{e}\frac{w_{e}(\lambda_{e}-\theta_{e})^{2}}{\nu_{e}},
\qquad \Gamma=\sup_{w\in\Delta}J(w).
\label{eq:game}
\end{equation}
The summand is the per-query Gaussian divergence, so $1/\Gamma$ is the rate of a
$\delta$-correct learner.

\leanbox{\lean{residual\_eq\_inner}, \lean{loss\_eq\_half\_range},
\lean{selectedRadius\_eq\_compressionLoss} in \lean{LeanPaperAModelBridgeV1};
\lean{energy\_eq\_sum\_gaussKL} in \lean{LeanPaperAStrictV1}}

\paragraph{The construction} Two arms, $A$ with cells $a_{0},a_{1}$ and $B$ with anchors
$b_{2},b_{3}$ and $2r$ middle cells grouped in $r$ pairs. The response is additive,
$F_{\lambda}(a,b)=\lambda_{A}(a)+\lambda_{B}(b)$. Anchors and arm-$A$ cells have variance $1$
and middle cells variance $M>0$. The world $\theta_{C}$ puts arm $A$ at $(-\frac12,\frac12)$,
the anchors at $(-\frac18,\frac18)$ and every middle cell at $0$. The class $\Theta_{0}$
requires each arm's cells to sum to zero.

\begin{theorem}[The construction is an observation model]
\label{thm:bridge}
$\lambda\in\Theta_{0}$ if and only if both uniform-product marginals recover their coordinates.
On $\Theta_{0}$, $L_{\lambda}(\{A\})=\frac12\mathrm{range}(\lambda_{B})$ and
$L_{\lambda}(\{B\})=\frac12\mathrm{range}(\lambda_{A})$. Retaining $A$ is $\varepsilon$-unsafe
if and only if $\mathrm{range}(\lambda_{B})>\mathrm{range}(\lambda_{A})+2\varepsilon$, and this
set is exactly the union over pairs $(p,q)$ of arm-$B$ cells of the polyhedra
$\{\lambda_{1}-\lambda_{0}+2\varepsilon<\lambda_{p}-\lambda_{q},\
\lambda_{0}-\lambda_{1}+2\varepsilon<\lambda_{p}-\lambda_{q}\}$.
\end{theorem}

\leanbox{\lean{centred\_iff\_marginals}, \lean{lossRetainA\_eq}, \lean{unsafeA\_iff\_range\_gap}
in \lean{LeanPaperACentredModelV1}; \lean{unsafe\_iff\_union\_of\_polyhedra} in
\lean{LeanPaperARegionV1}}

The unsafe condition is strict. Its closure contains ties, which are not unsafe. The game the
decision defines is therefore taken over the \emph{strict} alternatives
$\mathcal A_{\rm true}=\{\mathrm{range}(\lambda_{A})<\mathrm{range}(\lambda_{B})\}\cap\Theta_{0}$,
and every bound below is approached from inside that strict set.

\paragraph{The two games} The true game uses $\mathcal A_{\rm true}$. The frozen surrogate
keeps arm $A$'s cells and arm $B$'s anchors as the extrema and tests the one contrast
$(\lambda_{a_{1}}-\lambda_{a_{0}})-(\lambda_{b_{3}}-\lambda_{b_{2}})\le0$ over $\Theta_{0}$. The
two sets are not nested. The surrogate admits parameters that break the ordering, and it
misses alternatives that move a middle cell past an anchor.

\paragraph{A margin certificate is not a decision} Suppose the certificate condition implies
that $T$ is safe. Then the certificate's alternatives contain the decision's, so certifying is
never cheaper than deciding. On $F=ax_{1}+bx_{2}+cx_{1}x_{2}$ the decision is the sign of $a-b$,
while the usual margin certificate asks $a-b\ge2c$. At $a=3$, $b=1$, $\varepsilon=0.1$ the
certificate costs $19.4$ times the decision at $c=0.25$ and $1089$ times at $c=0.9$. It is
unattainable for $c\ge1$, where the decision is still settled. \NUM{}

\leanbox{\lean{certification\_costs\_at\_least\_as\_much} in
\lean{LeanPaperAOracleCertificateV1}}

\section{Separation after the allocation is optimised}
\label{sec:separation}

\begin{theorem}[Separation]
\label{thm:sep}
For every $r\ge1$ and $M>0$,
\[
\Gamma_{\rm fr}=\frac{9}{512},\qquad
\frac{9}{128(4+2rM)}\ \le\ \Gamma_{\rm true}\ \le\ \frac{1}{8Mr},
\]
and hence
\[
\frac{4}{4+2rM}\ \le\ \frac{\Gamma_{\rm true}}{\Gamma_{\rm fr}}\ \le\ \frac{64}{9rM}.
\]
The same bounds hold for the game over the strict alternatives. There, the Gaussian divergence
of the observation law is the energy of~\eqref{eq:game} for every centred candidate.
\end{theorem}

\leanbox{\lean{frozenGameC₀\_eq}, \lean{trueGameC₀\_le}, \lean{trueGameC₀\_ge},
\lean{centred\_separation} in \lean{LeanPaperACentredV1}; \lean{strict\_game\_bounds},
\lean{actualKLValue\_eq} in \lean{LeanPaperAEndgameV1}; \lean{gaussianKLIdentification\_holds}
in \lean{LeanPaperAGaussianKLV1}}

\emph{Why.} The upper bound is a pigeonhole. Some pair of middle cells carries at most $1/r$
of any design, and moving that pair apart by $\frac12+d$ reverses the decision. The move is a
strict alternative inside $\Theta_{0}$ for every $d>0$. The lower bound is the
variance-proportional design $w_{e}=\nu_{e}/(4+2rM)$, at which the energy is the plain squared
distance. Reversing the decision forces a contrast of at least $g=\frac34$ across four cells,
and Cauchy--Schwarz turns that into $\frac{9}{64}$ of squared distance. The same $g$ gives
$\Gamma_{\rm fr}=g^{2}/32$. Both infima are restricted to $\Theta_{0}$, and restricting an
infimum can only raise it. So neither bound is inherited from a larger set: each witness is
proved to lie in the class.

\section{The surrogate misdesigns}
\label{sec:starve}

A ratio of optimal values can diverge while both objectives select the same design. What
matters is the design the surrogate selects.

\begin{theorem}[Starvation]
\label{thm:starve}
For every allocation $w$, $J_{\rm fr}(w)\le\frac{9}{512}\,a(w)$, where $a(w)$ is the mass on
arm $A$ and the anchors. Consequently:
\begin{enumerate}
\item[(i)] every $w$ with $J_{\rm fr}(w)=\Gamma_{\rm fr}$ puts zero mass on the middle cells,
and $J_{\rm true}(w)=0$;
\item[(ii)] every $w$ with $J_{\rm fr}(w)\ge\Gamma_{\rm fr}-\varepsilon$ puts at most
$\frac{512}{9}\varepsilon$ on the middle cells, and $J_{\rm true}(w)\le\frac{64\varepsilon}{9rM}$.
\end{enumerate}
\end{theorem}

\leanbox{\lean{frozenValueC₀\_le}, \lean{starvation\_at\_optimum₀}, \lean{near\_optimal\_starvation₀},
\lean{near\_optimal\_trueValueC₀\_le} in \lean{LeanPaperACentredV1};
\lean{frozen\_optimum\_trueValueStrict\_zero} in \lean{LeanPaperAStrictV1}}

Solving the surrogate more accurately therefore drives the design towards one that cannot
distinguish the alternatives the compression turns on. Part~(ii) is what a learner running on
the surrogate actually has: an approximate optimum, which leaks in proportion to its own
suboptimality.

\section{The price of repair}
\label{sec:repair}

Write the \emph{relative sacrifice} of a design as
$\mathrm{sac}(w)=1-J_{\rm fr}(w)/\Gamma_{\rm fr}$. It is a diagnostic of a design, evaluated at
the true world, and it never enters a learner.

\subsection{At a design}

\begin{theorem}[The repair frontier]
\label{thm:frontier}
For every $r\ge1$, $M>0$ and $\beta\in[0,1]$, the best true information of a design with
$\mathrm{sac}(w)\le\beta$ satisfies
\[
\frac{9\beta}{128(4+2rM)}\ \le\ \sup_{\mathrm{sac}(w)\le\beta}J_{\rm true}(w)\ \le\
\frac{\beta}{8rM}.
\]
This holds on the strict game as well. Pointwise, every allocation has
$J_{\rm true}(w)\le\mathrm{sac}(w)/(8rM)$. The lower bound is attained by
$w_{\beta}=(1-\beta)w_{\rm unif}+\beta w_{\rm cert}$.
\end{theorem}

\leanbox{\lean{value\_concave}, \lean{frontier\_two\_sided} in \lean{LeanPaperAFrontierV1};
\lean{exchange\_strict}, \lean{strict\_frontier} in \lean{LeanPaperASacrificeV1}}

Each game's value is an infimum of functions linear in the design, hence concave in it, so
mixing two designs loses at most linearly. The upper side is Theorem~\ref{thm:starve}(ii) with
$\varepsilon=\beta\Gamma_{\rm fr}$. The exchange rate is linear in the sacrifice, and it
carries the same $\frac{1}{rM}$ as the separation.

\subsection{For a policy}

A design frontier does not say what a learner must give up. That requires every policy, its
own random sample sizes, and the correctness it must reach. Fix a randomised adaptive policy
with a finite horizon, whose answer event has probability at least $1-\delta$ at $\theta_{C}$
and at most $\delta$ at every strict alternative. Let $N_{e}$ be its expected number of queries
of cell $e$ at the truth, $N=\sum_{e}N_{e}$, and $\bar w=N_{\cdot}/N$ its average design. Put
$D=(1-\delta)\log(1/\delta)-\log2$.

\begin{theorem}[The sacrifice law]
\label{thm:sacrifice}
Every such policy queries each middle pair at least $8MD$ times in expectation, and the middle
cells at least $8rMD$ times in total. Each middle query is paid for in surrogate value: the
middle mass of any allocation is at most its relative sacrifice. Hence
\[
N\cdot\mathrm{sac}(\bar w)\ \ge\ 8\,r\,M\,D .
\]
The hypotheses are satisfiable. The ideal learner of Theorem~\ref{thm:ideal}, truncated at a
Markov horizon, satisfies them at level $2\delta$. For that learner
$N\,\mathrm{sac}(\bar w)\le N\le\mathbb{E}\tau_{\delta}$, and
$1/\Gamma_{\rm true}\le\frac{128(4+2rM)}{9}$.
\end{theorem}

\leanbox{\lean{pair\_demand}, \lean{middle\_demand}, \lean{middle\_le\_sacrifice},
\lean{sacrifice\_cost\_law}, \lean{upper\_constant}, \lean{constant\_gap} in
\lean{LeanPaperASacrificeV1}; \lean{ideal\_sacrifice}, \lean{sacrifice\_hypotheses\_satisfiable}
in \lean{LeanPaperASacrificeIdealV1}}

\emph{What this says.} To make the true decision, a policy must give up a \emph{budget} of
surrogate value, $8rM\log(1/\delta)$ query-units to first order. It is not a fraction. A policy
may spend the budget slowly, with a small sacrifice over a long run, or quickly, but it cannot
avoid it. The rate-optimal learner spends at most $\frac{128(4+2rM)}{9}\log(1/\delta)$ to first
order. So the budget is $\Theta(rM\log(1/\delta))$, and the two constants differ by
$\frac{32}{9}+\frac{64}{9rM}$.

\begin{theorem}[Spreading the budget]
\label{thm:dilute}
Let $\pi_{0}$ be a policy as above with $N_{0}$ expected real queries and average design
$\bar w_{0}$. Run it for its $n$ rounds, then query the four $A$ cells round-robin for $4k$ more
rounds and answer from the first $n$. For every $k$, this policy is $\delta$-correct with the same
answer event. It has $N=N_{0}+4k$ expected real queries, it obeys the sacrifice law, and it spends
no more budget than $\pi_{0}$:
\[
8rMD\ \le\ N\,\mathrm{sac}(\bar w)\ \le\ N_{0}\,\mathrm{sac}(\bar w_{0}),\qquad
\mathrm{sac}(\bar w)\ \le\ \frac{N_{0}}{N_{0}+4k}.
\]
\end{theorem}

\leanbox{\lean{sacrifice\_mix\_le}, \lean{padPol}, \lean{pad\_real}, \lean{lintegral\_cnt\_pad},
\lean{padCnt\_A}, \lean{dilution}, \lean{dilution\_small} in \lean{LeanPaperASacrificeDiluteV1}}

With $\pi_{0}$ the truncated ideal learner, $N_{0}\le\mathbb{E}\tau_{\delta}
=O(rM\log(1/\delta))$. So for every run length $N\ge N_{0}$, in steps of four, a
$\delta$-correct policy gives up a relative surrogate value of at most $N_{0}/N$. The budget
$N\,\mathrm{sac}(\bar w)$ stays between $8rMD$ and $N_{0}$. The padding is surrogate-optimal, and
concavity of the frozen value turns the mixture of designs into
$\mathrm{sac}\le\beta\,\mathrm{sac}(\bar w_{0})$ with $\beta=N_{0}/N$.

\emph{Why.} Each middle pair is moved by a strict alternative that only that pair's queries can
detect. The change of measure at the shift $\frac12+d$, with $d\to0$, gives the per-pair demand.
The surrogate's value is at most $\Gamma_{\rm fr}$ times the non-middle mass, which gives the
exchange.

\subsection{Forcing schedules}

A learner that follows the surrogate meets the demand for middle queries only through its
exploration. The sacrifice law then fixes its cost, whatever the schedule.

\begin{theorem}[The forcing law]
\label{thm:forcing}
For a policy as above, let $T$ be the number of real queries and $\mathrm{Mid}$ the number of
middle queries at the horizon. Suppose $\mathrm{Mid}\le\varphi(T)+B$ on every history, where
$\varphi$ is continuous and concave on $[0,\infty)$ and $\mathbb{E}B\le b$. Then
$\varphi(\mathbb{E}T)\ge8rMD-b$. For $\varphi(t)=ct^{a}$ with $0<a\le1$ and
$b\le 8rMD$,
\[
\mathbb{E}T\ \ge\ \Bigl(\frac{8rMD-b}{c}\Bigr)^{1/a}.
\]
\end{theorem}

\leanbox{\lean{forcing\_law}, \lean{power\_forcing\_law} in \lean{LeanPaperASacrificeV1}}

At $a=\frac12$, square-root forcing costs $\Omega(\log^{2}(1/\delta))$. At $a=1$, linear
forcing costs only $\Omega(\log(1/\delta))$, and the learner is no longer following the
surrogate. Whether a given learner satisfies the pathwise hypothesis is a property of that
learner. Section~\ref{sec:learners} proves it for the two surrogate learners.

\section{Learners}
\label{sec:learners}

\begin{theorem}[The ideal learner]
\label{thm:ideal}
There is a completely defined learner with the following steps. It starts warm, with one query
per cell. At round $s$ it selects an answer and a design by an exact maximiser of $J$ over a
grid of mesh $\frac{1}{s+1}$. It mixes in the uniform design with weight
$\frac{1}{\sqrt{s+1}}$ and tracks the result. It stops when some answer's generalised
likelihood ratio exceeds $\log(1/\delta)+\frac12\sum_{e}\log(1+N_{e}/a)$. On the construction,
for every $a>0$:
\begin{itemize}
\item it stops with a wrong answer with probability at most $\delta$;
\item for every $\eta\in(0,\Gamma_{\rm true})$ there is $A_{\eta}$ such that, for every
$\delta\in(0,1)$, $\mathbb{E}\tau_{\delta}\le\frac{\log(1/\delta)}{\Gamma_{\rm true}-\eta}+A_{\eta}$;
\item every finite-horizon policy that is $\delta$-correct against the strict alternatives
makes at least $\frac{(1-\delta)\log(1/\delta)-\log2}{\Gamma_{\rm true}}$ real queries in
expectation, and the stopped learner meets this bound at level $2\delta$.
\end{itemize}
Hence $\mathbb{E}\tau_{\delta}/\log(1/\delta)\to1/\Gamma_{\rm true}$, where $\Gamma_{\rm true}$ is
the Gaussian-divergence game over the strict alternatives.
\end{theorem}

\leanbox{\lean{ideal\_family\_finish}, \lean{ideal\_family\_correct}, \lean{ideal\_family\_lower},
\lean{GamT\_fam} in \lean{LeanPaperA7IdealFamilyV1}; \lean{ideal\_expect\_bound},
\lean{ideal\_upper\_rate} in \lean{LeanPaperA7IdealUpperV1}; \lean{family\_policy\_lower} in
\lean{LeanPaperAFamilyLowerV1}}

The world enters the learner only as the law of its observations. Correctness is Ville's
inequality at the learner's own statistic. The upper bound comes from box-free accumulation of
the statistic on the rounds where the confidence event holds, and that event fails at round
$s$ with probability at most $s^{-4}$. The lower bound couples the stopped learner, path by
path, with a policy that queries a fresh dummy after stopping. This learner is not the
implementation measured in Section~\ref{sec:evidence}: that one has finite precision and a
bracketing oracle, and its rate is only conditional (supplement).

\begin{theorem}[What following the surrogate costs]
\label{thm:surlearner}
Let $K=4+2r$. Consider the learner with the same warm start, forcing $\frac{1}{\sqrt{s+1}}$,
tracking and stop as the ideal learner, whose oracle returns an exact optimum of the frozen
surrogate. There are two cases:
\begin{itemize}
\item the surrogate is solved at the true world and held fixed;
\item the unrestricted surrogate is re-solved at the learner's current estimate every round.
The oracle must be exact at every estimate whose extrema are the four anchors and which
retains $A$. Elsewhere it is unconstrained.
\end{itemize}
In both cases there is $\delta_{0}>0$ with
\[
\mathbb{E}\tau_{\delta}\ \ge\ \Bigl(\frac{KM}{8}\Bigr)^{2}\log^{2}(1/\delta)
\qquad(\delta<\delta_{0}),
\]
so the ratio to the ideal learner is unbounded.
\end{theorem}

\leanbox{\lean{surrogate\_optimum\_cost}, \lean{surrogate\_vs\_ideal} in
\lean{LeanPaperA6SurrogateCostV1}; \lean{est\_cost\_log\_sq}, \lean{surrogate\_learner\_vs\_ideal}
in \lean{LeanPaperA6EstCostV1}; \lean{starvation\_at} in \lean{LeanPaperA6SurrogateAtV1}}

\emph{Why.} By starvation the oracle puts nothing on the middle cells, so forcing is the only
source of middle queries. The first case is an \emph{oracle-allocation} result: the true
extrema choose the starved cells. In the second case the oracle sees only the estimate. Forcing
gives every count order $\sqrt t$, so after a deterministic time every estimate lies within
$\frac1{16}$ of the truth, except on rounds where the learner's confidence event fails. Those
rounds number at most a constant in expectation at every horizon. Along every path the middle
count is then at most a multiple of $\sqrt{\tau}$ plus that count. This is the shape the
forcing law assumes at $a=\frac12$. The Lean proofs of this theorem establish it for these two
learners directly, with the stopped change of measure at level $2\delta$ (supplement).

\emph{Which surrogate.} The estimate-time learner solves the \emph{unrestricted} frozen game at
its estimate, without centring. At the truth this game has the same value, $\frac{9}{512}$, as
the centred game of Theorem~\ref{thm:sep}. Elsewhere the two differ. Shift both arm-$A$ cells
of the truth by $u$. The unrestricted surrogate of the uniform-anchor design stays at
$\frac{9}{512}$, while the centred one is at least $\frac{9}{512}+\frac{u^{2}}{4}$. The
theorems above are about the object they name.

\leanbox{\lean{frozenGameAt\_thetaC}, \lean{frozenAt\_unif\_shift}, \lean{centred\_unif\_shift\_ge},
\lean{two\_surrogates\_differ} in \lean{LeanPaperA6TwoSurrogatesV1}}

\section{Evidence}
\label{sec:evidence}

Every number below is produced by a script that checks itself against a statement proved in
Lean, and refuses to write output if a check fails. Oracle arms read the true world. They are
labelled as oracle-allocation diagnostics and are never competitors. Error rates are tested on
the right quantity: $\delta$-correctness bounds $\Pr(\text{stop and wrong})$ over launched runs,
not the error rate among stopped runs, and censoring is reported separately.

\begin{table}[t]
\caption{Evidence packages. Each script self-checks against a Lean statement.}
\label{tab:evidence}
\centering
\footnotesize
\resizebox{\columnwidth}{!}{%
\begin{tabular}{@{}llr@{}}
\toprule
Package & Script & Checks \\
\midrule
Certified separation, $r=1,2,4,8$ & \texttt{EXP\_paperA\_certified\_separation} & 6/6 \\
\quad its figure & \texttt{FIG\_paperA\_certified} & 3/3 \\
Allocation mechanism, one certificate & \texttt{EXP\_paperA\_mechanism} & 5/5 \\
\quad its figures & \texttt{FIG\_paperA\_mechanism}, \texttt{FIG\_paperA\_budget} & 7/7, 5/5 \\
Deployable rules, two-answer certificate & \texttt{EXP\_paperA\_deployable} & 6/6 \\
\quad its figure & \texttt{FIG\_paperA\_deployable} & 4/4 \\
\quad locked paired panel, $r\le4$, $M\in\{1,4\}$ & \texttt{EXP\_paperA\_deployable\_paired} & 5/5 \\
\quad its figure and paired table & \texttt{FIG\_paperA\_deployable\_paired} & 7/7 \\
Shipped learner, four-cell panel & \texttt{EXP\_paperA7\_panel} & 6/6 \\
\quad its figure & \texttt{FIG\_paperA7\_panel} & 4/4 \\
Error over launched runs & \texttt{ANALYSE\_paperA\_alpha\_launched} & 4/4 \\
Does the rule learn the answer? & \texttt{EXP\_paperA\_swap} & 7/7 \\
\bottomrule
\end{tabular}}
\end{table}

\begin{figure}[t]
\centering
\includegraphics[width=\columnwidth]{paperA_certified.pdf}
\caption{The separation, certified. A numerical true-game computation with centring and every
alternative region, certified in exact rational arithmetic (\CERT). The upper end is the dual of
the outer LP over exactly verified centred alternatives. The lower end is the minimum over all
regions of a Lagrangian dual value at a design with every weight positive.
$\Gamma_{\rm true}/\Gamma_{\rm fr}$ lies in $[0.8372,0.8374]$, $[0.7214,0.7216]$,
$[0.5664,0.5667]$ and $[0.3982,0.3982]$ at $r=1,2,4,8$ ($M=1$), inside the band of
Theorem~\ref{thm:sep}. $\Gamma_{\rm fr}=\frac{9}{512}$ is the theorem.}
\label{fig:acert}
\end{figure}

\begin{figure*}[t]
\centering
\includegraphics[width=0.92\textwidth]{paperA_mechanism.pdf}
\caption{The mechanism, from one stored panel ($M=1$, $\alpha=0.05$, 24 runs per rule and $r$,
budget $B=6\times10^{4}$); nothing was re-solved or re-sampled. Left: attained lower
certificates of the true value, up to the theorem's upper bound. Centre: the share of queries
on the middle cells. The surrogate optimum puts nothing there, and its runs reach the middles
only through forcing. Right: $\mathbb{E}[\min(\tau,B)]$. Every surrogate run is censored at $B$,
which does not show that it never stops. SCOPE: the stopping test certifies answer $A$ only,
so this is an allocation diagnostic under one certificate, not a $\delta$-correct learner
(\texttt{FIG\_paperA\_mechanism}, 7/7).}
\label{fig:amech}
\end{figure*}

\begin{figure*}[t]
\centering
\includegraphics[width=0.92\textwidth]{paperA_deployable.pdf}
\caption{Deployable rules under one two-answer certificate. Each rule is a tracking learner
with forcing $\min(1,1/\sqrt t)$, and its oracle sees only the estimate. \emph{Surrogate at
estimate} re-solves the unrestricted frozen surrogate at the estimate every round.
\emph{Geometry} ascends the certificate of the answer the estimate retains. Left, centre: the
certificate rate by budget on the base world (answer $A$) and the swapped world (answer $B$),
for $r=1,2,4,8$. Right: $\mathbb{E}[\min(\tau,B)]$, $B=6\times10^{4}$, $\alpha=0.05$, 24 runs per
cell. None of the 576 launched runs stopped with a wrong answer. The 71 censored runs are all
surrogate-at-estimate runs on the base world, and the other two rules certify every run. On the
swapped world the surrogate's frozen cells are the informative ones, and it certifies every
run. SCOPE: the certificate is a valid lower bound on the likelihood ratio, not the exact
infimum, and $\alpha=0.05$ lies far above the $\delta_{0}$ of Theorem~\ref{thm:surlearner}. The
panel shows these rules' behaviour; it does not test the $\log^{2}$ order
(\texttt{EXP\_paperA\_deployable}, 6/6; \texttt{FIG\_paperA\_deployable}, 4/4). The locked
paired panel of Figure~\ref{fig:apaired} repeats it at scale.}
\label{fig:adeploy}
\end{figure*}

\begin{figure*}[t]
\centering
\includegraphics[width=0.92\textwidth]{paperA_deployable_paired.pdf}
\caption{The locked panel. The same three rules, certificate and stop as
Figure~\ref{fig:adeploy}, with the configuration fixed and hashed before launch. There are
200 paired runs per cell: run $i$ of every rule in a cell uses the same seed. The grid is
$r\in\{1,2,4\}$, $M\in\{1,4\}$, both worlds, $B=6\times10^{4}$ and $\alpha=0.05$. Left, centre:
$\mathbb{E}[\min(\tau,B)]$ on the base and swapped worlds, with $M=1$ solid and $M=4$ dashed.
Right: the fraction certified by $B$, with exact 95\% intervals. Of 7200 launched runs, 6258
stopped, and none stopped with a wrong answer (0/2400 per rule). On the base world the
surrogate-at-estimate rule certifies 167, 41, 44, 2, 4 and 0 of 200 runs across the six
$(r,M)$ cells, while the other two rules certify every run. Its paired excess over geometry is
$1.9\times10^{4}$ to $5.4\times10^{4}$ rounds, with a bootstrap interval excluding $0$ in every
cell. Its median middle share is $0.05$, against $0.47$ for geometry. On the swapped world
every rule certifies every run, and the surrogate rule is the \emph{fastest} in every cell,
since there its frozen cells are the informative ones. SCOPE: $\min(\tau,B)$ is a restricted
mean, so where the surrogate rule is censored its mean and its excess are lower bounds. The
certificate is a lower bound and $\alpha$ lies far above $\delta_{0}$, so this measures these
rules and does not test the $\log^{2}$ order. All 31 cell summaries of an interrupted earlier
launch are reproduced exactly (\texttt{EXP\_paperA\_deployable\_paired}, 5/5;
\texttt{FIG\_paperA\_deployable\_paired}, 7/7).}
\label{fig:apaired}
\end{figure*}

\begin{figure*}[t]
\centering
\includegraphics[width=0.92\textwidth]{paperA7_panel.pdf}
\caption{The shipped learner on a four-cell panel: base world at $M=16$ and $M=64$, a smaller
gap, and the $A/B$-swapped world. 50 runs per cell and arm, $\alpha=0.05$. Left:
$\mathbb{E}[\min(\tau,B)]$, censored runs counted. Centre: the certified lower value of the true
game at the empirical design $N/\tau$, over the game's value. Offline truth evaluates it; it
never enters the learner. Right: selection. None of the 600 launched runs stopped with a wrong
answer; one learner run on the swapped world was censored. The panel exposes an algorithmic
weakness. The learner's kept set never narrowed to a single answer in any cell, so its
allocation follows the tie-break answer. On the swapped world that answer is the wrong one, the
learner selected the true answer in none of its refreshes, and it cost $3801$ queries against
$1109$ for the oracle-allocation diagnostic and $1679$ for uniform. Its stop was still correct.
SCOPE: cost estimation on a finite panel, not error certification
(\texttt{EXP\_paperA7\_panel}, 6/6; \texttt{FIG\_paperA7\_panel}, 4/4).}
\label{fig:apanel}
\end{figure*}

\section{Limitations}
\label{sec:limits}

\begin{enumerate}
\item[1)] \textbf{One construction.} Theorems~\ref{thm:sep}--\ref{thm:surlearner} are about the
centred construction. Two heteroscedastic families on which the surrogate misprices in opposite
directions are computed in the supplement. They are stated for a range oracle, and only the
overpricing direction is realised in the model.
\item[2)] \textbf{Spreading the budget, not reducing it.} Theorem~\ref{thm:dilute} spreads a
paid budget over any run length by padding after the decision is made. It does not show that
interleaving surrogate-optimal queries \emph{during} the run lowers the budget below $N_{0}$.
The lower bound $8rMD$ and the ideal learner's $N_{0}$ still differ by the constant of
Theorem~\ref{thm:sacrifice}.
\item[3)] \textbf{The implemented learner.} The learner of Figure~\ref{fig:apanel} has finite
precision and a bracketing oracle. Its rate is conditional on named accumulation and
confidence hypotheses (supplement), and its selection step did not narrow on the panel.
\item[4)] \textbf{Finite horizon, finite answers.} The sacrifice and forcing laws are stated for
finite-horizon policies. The ideal learner's lower bound reaches unbounded stopping times
through truncation at a Markov horizon.
\item[5)] \textbf{Certificates.} The regional KKT checks behind the two families are
floating-point. The separation of Figure~\ref{fig:acert} is certified exactly.
\end{enumerate}

\section{Conclusion}

A frozen surrogate is not a small error in a correct model. After the allocation is optimised
it overstates the available information by a factor $\Theta(rM)$. Its optima collect nothing
about the decision, and a learner that follows it pays a whole extra power of $\log(1/\delta)$.
Undoing it has a definite price. At a design, true information comes back linearly in the
sacrifice, at the separation's own rate. For a policy, the true decision requires a budget of
$\Theta(rM\log(1/\delta))$ query-units of surrogate value, which no schedule avoids and the
rate-optimal learner attains in order.

\section*{Companion material}
The supplement gives the proof architecture of every theorem, the supporting Lean declarations,
the conditional status of the implemented learner's rate, the heteroscedastic families, the
repair criterion and the confidence-sequence machinery, with each superseded statement marked.
The claim map records each headline's model, declaration, scope and evidence.

\begin{thebibliography}{00}

\bibitem{degenne2019} R. Degenne and W. M. Koolen, ``Pure exploration with multiple correct
answers,'' in \emph{Advances in Neural Information Processing Systems 32}, 2019,
pp.~14591--14600.

\bibitem{wangyou2026} Q. Wang and W. You, ``Ranking-and-selection with multiple correct
answers and non-answerable estimates,'' arXiv:2606.21889, 2026.

\bibitem{bennouna2026} O. Bennouna, A. Bennouna, S. Amin, and A. Ozdaglar, ``Data
informativeness in linear optimization under uncertainty,'' arXiv:2602.15365, 2026.

\bibitem{schutte2026} N. Schutte et al., ``Sufficient decision proxies for
decision-focused learning,'' arXiv:2505.03953, 2026.

\bibitem{overstall2022} A. M. Overstall, ``Properties of Fisher information gain for
Bayesian design of experiments,'' Univ. of Southampton, 2022. [Online]. Available:
https://eprints.soton.ac.uk/452554/

\bibitem{filstroff2021} L. Filstroff, I. Sundin, P. Mikkola, A. Tiulpin, J. Kylm\"{a}oja,
and S. Kaski, ``Targeted active learning for Bayesian decision-making,'' arXiv:2106.04193,
2021.

\bibitem{tang2026} Y. Tang, S. Sloman, and S. Kaski, ``Error amplification under model
misspecification,'' in \emph{Proc. Machine Learning Research}, vol.~300, 2026.

\bibitem{wan2026} Z. Wan, Y. Liu, P. Grigas, and Z.-J. M. Shen, ``Decision-focused
sequential experimental design,'' arXiv:2602.05340, 2026.

\bibitem{mathlib2020} The mathlib Community, ``The Lean mathematical library,'' in
\emph{Proc. 9th ACM SIGPLAN Int. Conf. Certified Programs and Proofs (CPP)}, 2020,
pp.~367--381.

\bibitem{demoura2021} L. de Moura and S. Ullrich, ``The Lean 4 theorem prover and
programming language,'' in \emph{Proc. 28th Int. Conf. Automated Deduction (CADE)}, 2021,
pp.~625--635.

\end{thebibliography}

\end{document}
